\section{Conclusion}
We classified cardiac pathologies from ED and ES segmentations, separating strictly the data used to select models from the data used to score them. On 21 volumetric biomarkers, selecting hyper-parameters on a single split gives accuracies of 0.91 to 0.95, whereas nested cross-validation gives 0.90 to 0.94 for the same models, none of which stands out; ensemble weights chosen on the evaluation folds are noise. The missing LV cavity of the test patients is recovered almost exactly by filling the myocardium, and a convex-hull fallback fixes the open basal slices; whether training uses the true or the reconstructed cavity makes no measurable difference. The step that improves classification is describing the myocardial wall: thickness features raise the accuracy by 3 to 6 points, and an average of an MLP ensemble and a linear SVM reaches $0.966 \pm 0.011$. Its remaining errors are infarcted hearts that look dilated and the converse. Future work should evaluate the pipeline on automatic segmentations of the test images, add features of the whole cardiac cycle such as regional wall motion, and score the predictions on held-out labels.
